\documentclass[11pt,a4paper]{article}
\usepackage[a4paper,margin=2.3cm]{geometry}
\usepackage{amsmath,amssymb}
\usepackage{booktabs,array,tabularx}
\usepackage[group-separator={,},group-minimum-digits=4,exponent-product=\cdot]{siunitx}
\usepackage[hidelinks]{hyperref}
\usepackage{parskip}
\sisetup{per-mode=symbol,group-digits=integer,range-phrase={--},range-units=single,list-separator={, },list-final-separator={, and },list-pair-separator={ and }}
\sloppy
\title{Structural Integrity of an Ice Shell\\Concentric Torus Habitat in Arcadia Planitia, Mars}
\author{Concept-level calculation}
\date{November 2026}
\begin{document}
\maketitle

\begin{abstract}
This note checks whether a self-supporting ice vault (one ring of a torus habitat) can stand for spans of \SIlist{16;24;32;40}{\metre}, and how much heat from the habitat affects the ice. Compressive stress in the ice is small (\SIrange{46}{115}{\kilo\pascal}), but creep-driven buckling sets the shell thickness: \SI{0.6}{\metre} for \SIrange{16}{24}{\metre} spans, \SI{1.25}{\metre} for \SI{32}{\metre} and \SI{3.4}{\metre} for \SI{40}{\metre}. The governing variable is not span but ice temperature: the shell must stay near \SI{-60}{\celsius}, with insulation on the inner face ($R \geq \SI{4}{\metre\squared\kelvin\per\watt}$) and a reflective cover on the outer face. All numbers are concept-level estimates.
\end{abstract}

\section{Summary of results}
\begin{table}[h]\centering\small
\caption{Summary by span. FS = buckling safety factor after 30 years of creep. ``Damaged'' = reflective cover lost (\S\ref{sec:bad}).}
\begin{tabular}{rrrrrrrr}\toprule
$L$ & $h$ & $t$ & $\sigma_{\max}$ & FS & FS & Arch mass & Tie force \\
(\si{\metre}) & (\si{\metre}) & (\si{\metre}) & (\si{\kilo\pascal}) & design & damaged & (\si{\tonne}) & (\si{\kilo\newton}) \\\midrule
\num{16} & \num{9.0} & \num{0.60} & \num{46} & \num{23.6} & \num{2.93} & \num{35} & \num{23} \\
\num{24} & \num{13.5} & \num{0.60} & \num{69} & \num{3.3} & \num{0.39} & \num{53} & \num{35} \\
\num{32} & \num{18.0} & \num{1.25} & \num{92} & \num{2.6} & \num{0.31} & \num{146} & \num{97} \\
\num{40} & \num{22.5} & \num{3.40} & \num{115} & \num{2.5} & \num{0.37} & \num{498} & \num{330} \\
\bottomrule\end{tabular}
\end{table}
Arch mass and floor-tie force are per \SI{2.5}{\metre}-wide segment.

\section{Data and assumptions}
\begin{table}[h]\centering\small
\begin{tabularx}{\linewidth}{lrX}\toprule
Quantity & Value & Note \\\midrule
Mars gravity $g$ & \SI{3.71}{\metre\per\second\squared} & \\
Ice density $\rho$ & \SI{917}{\kilo\gram\per\metre\cubed} & bubble-free \\
Ice elastic modulus $E_0$ & \SI{9}{\giga\pascal} & short-term \\
Latent heat $L_f$ & \SI{334}{\kilo\joule\per\kilo\gram} & \\
Specific heat $c_p$ & \SI{1700}{\joule\per\kilo\gram\per\kelvin} & near \SI{-60}{\celsius} \\
Conductivity $k(T)$ & $9.828\,e^{-0.0057T}$ \si{\watt\per\metre\per\kelvin} & $T$ in K; $k(\SI{-60}{\celsius})=\,$\num{2.92} \si{\watt\per\metre\per\kelvin} \\
Habitat air $T_{in}$ & \SI{20}{\celsius} & \\
Mean outer-face temperature $T_o$ & \SI{-60}{\celsius} & with reflective cover; seasonal swing $\pm\SI{15}{\celsius}$ \\
Insulation, module to ice $R_{ins}$ & \SI{4}{\metre\squared\kelvin\per\watt} & e.g.\ aerogel blanket $\approx\SI{6}{\centi\metre}$ \\
Rise/span ratio $h/L$ & \num{0.5625} & same as the \SI{16}{\metre}~$\times$~\SI{9}{\metre} arch \\
Design life $t_d$ & 30 years & \\
Minimum thickness (shielding) & \SI{0.6}{\metre} & $\approx\SI{55}{\gram\per\centi\metre\squared}$ \\
\bottomrule\end{tabularx}
\end{table}
The \SI{56.5}{\kilo\pascal} habitat pressure is held by inflatable modules, not by the ice. Between module and shell is a gap of low-pressure Mars gas, so the ice carries only its own weight.

\section{Arch geometry}
The arch is an inverted \emph{catenary}. Under uniform self-weight along the arch, this shape puts the whole section in pure compression, with no bending:
\begin{equation}
y(x) = h - a\left[\cosh\frac{x}{a}-1\right], \qquad -\tfrac{L}{2}\le x\le \tfrac{L}{2}.
\end{equation}
The parameter $a$ follows from the feet sitting on the ground:
\begin{equation}
a\left[\cosh\frac{L}{2a}-1\right]=h .
\end{equation}
Arc length:
\begin{equation}
s = 2a\sinh\frac{L}{2a}.
\end{equation}

\section{Loads}
The main load is self-weight per unit shell area:
\begin{equation}
w=\rho g t .
\end{equation}
Other loads are very small:
\begin{itemize}
\item \textbf{Wind.} Dynamic pressure $q=\tfrac12\rho_{atm}v^2 \approx \tfrac12(0.02)(30)^2 \approx \SI{9}{\pascal}$, under \SI{1}{\percent} of the weight of a \SI{0.6}{\metre} shell ($\approx\SI{2.0}{\kilo\pascal}$). Neglected.
\item \textbf{Habitat pressure.} Carried by the inflatable modules; it never reaches the ice.
\item \textbf{Marsquakes.} Very weak compared with earthquakes. Neglected at concept stage.
\end{itemize}

\section{Compressive stress}
For a catenary, the horizontal force is constant, $H = w\,a$, and the axial force at height $y$ is $N = w\,(a + h - y)$. It peaks at the feet ($y=0$). Dividing by thickness $t$, the thickness cancels:
\begin{equation}
\boxed{\;\sigma_{\max} = \rho\, g\,(a+h)\;}
\end{equation}
So the compressive stress depends only on arch size, not thickness. The crushing strength of cold ice is taken as \SI{5}{\mega\pascal} (conservative).

\begin{table}[h]\centering\small
\caption{Geometry and compressive stress.}
\begin{tabular}{rrrrrr}\toprule
$L$ (\si{\metre}) & $h$ (\si{\metre}) & $a$ (\si{\metre}) & $s$ (\si{\metre}) & $\sigma_{\max}$ (\si{\kilo\pascal}) & FS crushing \\\midrule
\num{16} & \num{9.0} & \num{4.56} & \num{25.5} & \num{46.1} & \num{108} \\
\num{24} & \num{13.5} & \num{6.85} & \num{38.3} & \num{69.2} & \num{72} \\
\num{32} & \num{18.0} & \num{9.13} & \num{51.1} & \num{92.3} & \num{54} \\
\num{40} & \num{22.5} & \num{11.41} & \num{63.9} & \num{115.4} & \num{43} \\
\bottomrule\end{tabular}
\end{table}
The ice will not crush: the safety factor exceeds 40 for every span.

\section{Ice creep}
Ice under sustained stress flows slowly. We use Glen's law ($n=3$):
\begin{equation}
\dot\varepsilon_{ij}=A(T)\,\tau_e^{\,n-1}\,\tau_{ij}.
\end{equation}
For uniaxial compression $\sigma$: $\tau_{11}=\tfrac23\sigma$ and $\tau_e=\sigma/\sqrt3$, so
\begin{equation}
\dot\varepsilon=\tfrac{2}{9}\,A(T)\,\sigma^{3}.
\end{equation}
The rate factor $A$ depends strongly on temperature (Arrhenius; Cuffey \& Paterson 2010):
\begin{equation}
A(T)=A^\ast\exp\!\left[-\frac{Q_c}{R}\left(\frac1T-\frac1{T^\ast}\right)\right],\quad
A^\ast=\SI{3.5e-25}{\per\second\per\pascal\cubed},\; T^\ast=\SI{263}{\kelvin},\; Q_c=\SI{60}{\kilo\joule\per\mole}.
\end{equation}
Values: $A(\SI{-60}{\celsius})=$ \num{5.63e-28}, $A(\SI{-40}{\celsius})=$ \num{1.03e-26}, $A(\SI{-20}{\celsius})=$ \num{1.18e-25} \si{\per\second\per\pascal\cubed}. Every \SI{20}{\celsius} of warming makes ice creep roughly 10--20 times faster.

Total creep strain over the design life, and the effective modulus used for buckling:
\begin{equation}
\varepsilon_c=\tfrac29\,\bar A\,\sigma_{\max}^3\,t_d,\qquad
E_{\text{eff}}=\frac{\sigma_{\max}}{\sigma_{\max}/E_0+\varepsilon_c}.
\end{equation}
$\bar A$ is $A$ averaged through the shell thickness (temperature varies linearly from inner to outer face) and over one year (sinusoidal seasonal swing). Because $A$ rises exponentially with temperature, the warmest part dominates.

\section{Buckling}
A thin arch can snap out of shape and collapse even at low stress. For a two-hinged parabolic/catenary arch under uniform load (Timoshenko \& Gere 1961), the critical load per metre of width is
\begin{equation}
q_{cr}=\gamma\,\kappa\,\frac{E_{\text{eff}}\,I}{L^3},\qquad I=\frac{t^3}{12},
\end{equation}
with $\gamma\approx36$ for $h/L\approx0.56$ and an imperfection factor $\kappa=0.5$. The self-weight load per horizontal metre is $q=\rho g t\,s/L$. The safety factor is
\begin{equation}
\mathrm{FS}=\frac{q_{cr}}{q}=\frac{\gamma\,\kappa\,E_{\text{eff}}\,t^{2}}{12\,\rho\, g\, s\, L^{2}}
\quad\Longrightarrow\quad
t_{\text{req}}=L\sqrt{\frac{12\,\mathrm{FS}\,\rho\, g\, s}{\gamma\,\kappa\,E_{\text{eff}}\,L}} .
\end{equation}
Target $\mathrm{FS}\ge2.5$. The calculation is iterated because $E_{\text{eff}}$ depends on temperature, and the inner-face temperature depends on thickness (\S\ref{sec:thermal}). Design thickness is rounded up to \SI{5}{\centi\metre} and is never below \SI{0.6}{\metre} (radiation shielding).

\begin{table}[h]\centering\small
\caption{Required thickness against buckling (design case: $T_o=\SI{-60}{\celsius}\pm\SI{15}{\celsius}$, $R_{ins}=4$).}
\begin{tabular}{rrrrrr}\toprule
$L$ (\si{\metre}) & $t_{\text{req}}$ (\si{\metre}) & $t$ design (\si{\metre}) & $\varepsilon_c$ 30 yr (\si{\percent}) & $E_{\text{eff}}$ (\si{\giga\pascal}) & FS \\\midrule
\num{16} & \num{0.18} & \num{0.60} & \num{0.004} & \num{0.97} & \num{23.65} \\
\num{24} & \num{0.51} & \num{0.60} & \num{0.014} & \num{0.46} & \num{3.31} \\
\num{32} & \num{1.22} & \num{1.25} & \num{0.046} & \num{0.20} & \num{2.58} \\
\num{40} & \num{3.39} & \num{3.40} & \num{0.229} & \num{0.05} & \num{2.50} \\
\bottomrule\end{tabular}
\end{table}
Required thickness grows much faster than span: from \SI{24}{\metre} to \SI{40}{\metre} it rises more than fivefold. Two reasons: a larger arch loads the ice harder, so it creeps faster ($E_{\text{eff}}$ drops), and buckling stiffness falls with $L^3$. Stiffening ribs at segment joints would reduce this, because $I$ grows with the cube of section depth.

\section{Thermal behaviour and inner-surface melting}\label{sec:thermal}
\subsection{Heat-flow model}
Heat flows from the habitat air, through the insulation, through the ice, and is radiated to the sky from the outer face. Steady state, one dimension:
\begin{equation}
q=\frac{T_{in}-T_o}{R_{ins}+t/k},\qquad
T_i=T_o+q\,\frac{t}{k},
\end{equation}
where $q$ is the heat flux (\si{\watt\per\metre\squared}) and $T_i$ the inner-face ice temperature. Total heat loss for ring 1 uses the inner area $A_{in}=2\pi R_1 s$ with $R_1 = \SI{20}{\metre}+L/2$ (\O\SI{40}{\metre} courtyard). This is an upper bound, since the modules are smaller than the shell.

\subsection{Melting criterion}
The inner face reaches \SI{0}{\celsius} if the flux into the ice exceeds
\begin{equation}
q_{\text{melt}}=\frac{k\,(0-T_o)}{t}.
\end{equation}
If a local source (lamps, machinery, a person leaning on it, a leak of warm air) delivers more than $q_{\text{melt}}$, ice starts to melt at a rate
\begin{equation}
v=\frac{q-q_{\text{melt}}}{\rho\,L_f}.
\end{equation}
Example: an excess of \SI{150}{\watt\per\metre\squared} melts $\approx$ \num{4.2} \si{\centi\metre\per\day}. The time for heat to pass through the full thickness is
\begin{equation}
\tau=\frac{t^2}{\alpha},\qquad \alpha=\frac{k}{\rho c_p}=\,\num{1.87e-6}\;\si{\metre\squared\per\second}.
\end{equation}

\begin{table}[h]\centering\small
\caption{Thermal results, design case ($T_o=\SI{-60}{\celsius}$, $R_{ins}=\SI{4}{\metre\squared\kelvin\per\watt}$).}
\begin{tabular}{rrrrrrr}\toprule
$L$ (\si{\metre}) & $t$ (\si{\metre}) & $q$ (\si{\watt\per\metre\squared}) & $T_i$ (\si{\celsius}) & $q_{\text{melt}}$ (\si{\watt\per\metre\squared}) & $\tau$ (days) & Ring 1 heat loss (\si{\kilo\watt}) \\\midrule
\num{16} & \num{0.60} & \num{19.0} & \num{-56.0} & \num{288} & \num{2.3} & \num{85} \\
\num{24} & \num{0.60} & \num{19.0} & \num{-56.0} & \num{288} & \num{2.3} & \num{146} \\
\num{32} & \num{1.25} & \num{18.0} & \num{-52.1} & \num{137} & \num{9.9} & \num{208} \\
\num{40} & \num{3.40} & \num{15.3} & \num{-41.2} & \num{49} & \num{75.5} & \num{246} \\
\bottomrule\end{tabular}
\end{table}

With $R=4$ insulation the inner face stays at \SIrange{-56}{-41}{\celsius}, far from melting. Thicker shells run slightly warmer inside, because the ice itself adds thermal resistance. The values of $\tau$ (days to weeks) mean temperature sensors have time to catch a hot spot before any ice melts.

\subsection{Insulation and structural safety}
Heat from inside does not only risk melting; it warms the ice and speeds up creep. For a \SI{32}{\metre} span with $t=\SI{1.25}{\metre}$:
\begin{table}[h]\centering\small
\caption{Sensitivity to insulation (\SI{32}{\metre} span).}
\begin{tabular}{rrrrr}\toprule
$R_{ins}$ (\si{\metre\squared\kelvin\per\watt}) & $q$ (\si{\watt\per\metre\squared}) & $T_i$ (\si{\celsius}) & FS buckling & Heat loss (\si{\kilo\watt}) \\\midrule
\num{0.25} & \num{106.8} & \num{-6.7} & \num{0.05} & \num{1234} \\
\num{0.50} & \num{81.7} & \num{-20.8} & \num{0.17} & \num{944} \\
\num{1.00} & \num{54.8} & \num{-34.8} & \num{0.58} & \num{633} \\
\num{2.00} & \num{32.7} & \num{-45.4} & \num{1.48} & \num{378} \\
\num{4.00} & \num{18.0} & \num{-52.1} & \num{2.58} & \num{208} \\
\num{8.00} & \num{9.5} & \num{-55.9} & \num{3.45} & \num{110} \\
\bottomrule\end{tabular}
\end{table}
With too little insulation ($R<2$), the inner face warms to \SIrange{-35}{-7}{\celsius} and the buckling safety factor drops below 1: \textbf{the arch slowly sags until it collapses}, long before any ice melts. Creep, not melting, is the governing criterion. Thick insulation also cuts the habitat's heating demand.

\subsection{Daily and seasonal temperature swings}\label{sec:bad}
The penetration depth of a temperature wave is $d=\sqrt{\alpha P/\pi}$:
\begin{itemize}
\item daily (1 sol): $d\approx$ \num{0.23} \si{\metre}, only the outer skin;
\item seasonal (1 Mars year): $d\approx$ \num{6.0} \si{\metre}, deeper than the shell. The whole shell warms in summer.
\end{itemize}
Damaged case: the reflective cover is lost, so the mean outer-face temperature rises to \SI{-55}{\celsius} with a $\pm\SI{30}{\celsius}$ swing (summer peak \SI{-25}{\celsius}) for the whole design life.
\begin{table}[h]\centering\small
\caption{Damaged case (reflective cover lost).}
\begin{tabular}{rrrrrrr}\toprule
$L$ (\si{\metre}) & $t$ (\si{\metre}) & $T_i$ (\si{\celsius}) & $\varepsilon_c$ (\si{\percent}) & $E_{\text{eff}}$ (\si{\giga\pascal}) & FS & $t$ required (\si{\metre}) \\\midrule
\num{16} & \num{0.60} & \num{-51.2} & \num{0.04} & \num{0.121} & \num{2.93} & \num{0.55} \\
\num{24} & \num{0.60} & \num{-51.2} & \num{0.13} & \num{0.054} & \num{0.39} & \num{1.96} \\
\num{32} & \num{1.25} & \num{-47.4} & \num{0.39} & \num{0.024} & \num{0.31} & $>8$ \\
\num{40} & \num{3.40} & \num{-37.0} & \num{1.57} & \num{0.007} & \num{0.37} & $>8$ \\
\bottomrule\end{tabular}
\end{table}
Only the \SI{16}{\metre} span stays safe. Larger spans require a maintained reflective cover, temperature monitoring, and a repair plan.

\section{Thrust at the arch feet}
The horizontal force at the feet per metre of width, $H=\rho g t\, a$, is carried by basalt floor ties between the inner and outer feet. For a \SI{32}{\metre} span, the force per \SI{2.5}{\metre} segment is $\approx$ \SI{97}{\kilo\newton}. With basalt-bar tensile strength $\approx\SI{1000}{\mega\pascal}$ and FS 3, the required area is $\approx\SI{300}{\milli\metre\squared}$: about one \O\SI{20}{\milli\metre} bar per segment.

\section{Conclusions and design rules}
\begin{enumerate}
\item Compressive stress is not a problem for any span studied.
\item Thickness is set by creep-driven buckling: \SI{0.6}{\metre} (\SIrange{16}{24}{\metre}), \SI{1.25}{\metre} (\SI{32}{\metre}), \SI{3.4}{\metre} (\SI{40}{\metre}). A \SI{40}{\metre} span is impractical without ribs.
\item Ice temperature is the key design variable. Keep the inner face $\le\SI{-50}{\celsius}$ with insulation $R\ge\SI{4}{\metre\squared\kelvin\per\watt}$; keep the outer face near \SI{-60}{\celsius} with a reflective cover.
\item Inner-surface melting only occurs where local flux exceeds $q_{\text{melt}}$ (\SIrange{50}{290}{\watt\per\metre\squared}). Heat sources must not touch or shine directly on the ice; an air gap and insulation blanket are mandatory.
\item Fit thermistors and strain gauges in every segment. The thermal response time of several days leaves time to act.
\item Spans of 24--32 m are feasible for public spaces (sports hall, pool) if rule 3 is met.
\end{enumerate}

\section{Limitations}
The buckling model uses an effective-modulus method and two-hinged arch coefficients, not finite-element analysis. Each arch is treated as an independent strip; shell action between segments and ring curvature are ignored, so results are conservative. Creep parameters come from terrestrial glacier ice; ice frozen in bags may differ. Arcadia surface temperatures and the effect of the reflective cover should be checked with a climate model (Mars Climate Database). Sublimation from the outer face has not been calculated.

\section*{References}
\begin{itemize}
\item Cuffey, K.\,M. \& Paterson, W.\,S.\,B. (2010). \emph{The Physics of Glaciers}, 4th ed. Elsevier. (Glen's law, $A(T)$)
\item Glen, J.\,W. (1955). The creep of polycrystalline ice. \emph{Proc.\ R.\ Soc.\ A} 228, 519--538.
\item Petrenko, V.\,F. \& Whitworth, R.\,W. (1999). \emph{Physics of Ice}. Oxford University Press.
\item Timoshenko, S.\,P. \& Gere, J.\,M. (1961). \emph{Theory of Elastic Stability}, 2nd ed. McGraw-Hill.
\item Bramson, A.\,M. et al.\ (2015). Widespread excess ice in Arcadia Planitia, Mars. \emph{Geophys.\ Res.\ Lett.} 42.
\end{itemize}
\end{document}
